% ==================
% 2.2 Frontend
% ==================

\section{Frontend}
\setauthor{\secondauthor}

Das Frontend bildet die visuelle und interaktive Ebene der Anwendung. Es ist dafür verantwortlich, komplexe Datenstrukturen benutzerfreundlich aufzubereiten und eine performante Bedienung zu ermöglichen. Im Fokus steht hierbei die Umsetzung einer modernen Oberfläche, die durch kurze Reaktionszeiten und eine klare Benutzerführung überzeugt.

Dabei muss das Frontend folgende projektspezifische Kernaufgaben erfüllen:
\begin{itemize}
\item \textbf{Dynamische Datenvisualisierung:} Aufbereitung von Lernfortschritten und Fehleranalysen in Form von interaktiven Statistiken, um den Lernerfolg visuell darzustellen.
\item \textbf{Zustandsmanagement:} Effiziente clientseitige Verwaltung des Anwendungszustands, um einen flüssigen Wechsel zwischen verschiedenen Ansichten und Datenpools ohne spürbare Ladezeiten zu garantieren.
\item \textbf{Schnittstellen-Konsumation:} Integration der REST-API zur asynchronen Datenverarbeitung, wobei die Kommunikation mit dem Server so optimiert wird, dass eine hohe Reaktivität der Anwendung erhalten bleibt.
\end{itemize}

Die technologische Wahl des Frameworks ist entscheidend für die Performance und die langfristige Wartbarkeit der Benutzeroberfläche. Im folgenden Vergleich werden die infrage kommenden Technologien gegenübergestellt.

\vspace{1cm}
\textbf{Analyse technischer Anforderungen}
\setauthor{\secondauthor}

Die Auswahl des Frontend-Stacks ist eine strategische Entscheidung, die bestimmt, wie effizient die Benutzeroberfläche gerendert wird und wie strukturiert der Code aufgebaut ist. Dabei wurden vor allem die Wartbarkeit und die Tooling-Unterstützung bewertet.

Folgende Punkte wurden dabei kritisch untersucht:
\begin{itemize}
\item \textbf{Rendering-Verfahren und DOM-Manipulation:}
Untersuchung, wie das Framework Änderungen am UI vornimmt (z. B. Virtual DOM). Ziel ist eine hohe Performance, damit die Anwendung auch bei einer großen Anzahl an darzustellenden Elementen flüssig reagiert.
\item \textbf{Ökosystem und Community-Support:}
Bewertung der Verfügbarkeit von etablierten Bibliotheken (z. B. für Routing oder UI-Komponenten). 
\item \textbf{Entwickler-Produktivität und Tooling:}
Analyse der verfügbaren Entwicklerwerkzeuge. Es wurde geprüft, wie gut Debugging-Möglichkeiten und Hot-Reload-Funktionen unterstützt werden, um den Entwicklungsprozess zu beschleunigen.
\item \textbf{Styling und Design-System-Integration:}
Vergleich der Möglichkeiten zur Gestaltung der Oberfläche. Hierbei steht die einfache Integration von modernen CSS-Lösungen im Vordergrund, um ein konsistentes Design-System umzusetzen.
\item \textbf{Lernkurve und Syntax:}
Bewertung der Komplexität des Frameworks. Da die Effizienz in der Umsetzung entscheidend ist, wurde untersucht, wie schnell die Konzepte des Frameworks sicher beherrscht werden können.
\end{itemize}


% ==================
% 2.2.1 Angular
% ==================
\subsection{Angular}
\setauthor{\secondauthor}


\begin{figure}[H]
\centering
\fbox{

\includegraphics[width=0.45\textwidth]{Angular-js-logo.png}
}
\caption{Logo von Angular}
\label{fig:angular}
\end{figure}


\textbf{Historischer Kontext: Von AngularJS zum modernen Framework}


Um die heutige Rolle von Angular zu verstehen, lohnt sich ein Blick auf die Webentwicklung um das Jahr 2010. Zu dieser Zeit dominierten entweder rein statische Webseiten oder dynamische Effekte, die mit Bibliotheken wie jQuery manuell zusammengefügt wurden \cite{AngularHistoryOnline}. In dieser Phase veröffentlichte Google 2010 AngularJS (Version 1.x), um Ordnung und Struktur in die clientseitige Anwendungsentwicklung zu bringen. Das zentrale Konzept war das sogenannte Two-Way-Data-Binding, bei dem Änderungen im Modell automatisch die Benutzeroberfläche aktualisieren, und umgekehrt. Dadurch wurden erstmals komplexe Single-Page-Applications (SPAs) möglich, also Webanwendungen, die innerhalb einer einzigen HTML-Seite laufen und Inhalte dynamisch nachladen, ohne die Seite vollständig zu aktualisieren \cite{AngOverview}.


Mit wachsender Komplexität von Webanwendungen und steigenden Anforderungen in Unternehmen stieß AngularJS jedoch an seine Grenzen. Das datengetriebene Aktualisierungsmodell führte bei großen Datenmengen zu Performance-Problemen, und die fehlende Typisierung in JavaScript erschwerte die Wartung in großen Teams \cite{AngMigration}. Google entschied sich daher 2016 zu einem vollständigen Neustart: Angular (ab Version 2). Der Code wurde komplett in TypeScript neu entwickelt, um statische Typisierung und bessere Tool-Unterstützung zu bieten \cite{AngularNG2}. Zudem ersetzte man das ineffiziente Zwei-Wege-Datenbinding durch ein performanteres, unidirektionales Change-Detection-System \cite{AngularHistory}.


\textbf{Architektur: Ein integriertes Framework}


Im Unterschied zu spezialisierten Bibliotheken wie React versteht sich Angular als umfassendes Framework („batteries included''). Es bringt wesentliche Funktionen wie Routing, Formularvalidierung und HTTP-Kommunikation bereits im Kern mit \cite{AngularFrameworkPast}. Für größere Projekte bedeutet das: konsistente Strukturen, klare Verantwortlichkeiten und weniger Abhängigkeiten von Drittbibliotheken.


Angular organisiert jede Benutzeroberfläche in eine feste Dreiteilung aus TypeScript-Logik, HTML-Template und CSS-Styling. Diese Struktur fördert einheitlichen Code, erleichtert Wiederverwendung und skaliert gut, wenn mehrere Teams parallel entwickeln. Gleichzeitig erlaubt Angular komponentenspezifische Styles, etwa durch ViewEncapsulation, wodurch individuelle Gestaltung innerhalb eines klaren architektonischen Rahmens möglich bleibt.


\vspace{0.7cm}
\textbf{Technisches Herzstück: RxJS und das reaktive State-Management}


Ein wesentliches Kriterium für die Lernplattform ist die effiziente Verwaltung des Anwendungszustands (State-Management), um flüssige Übergänge zwischen Fragenkatalogen zu ermöglichen. Angular setzt hierbei konsequent auf \textit{RxJS} (Reactive Extensions) \cite{AngularSignalsGuide}.


\begin{figure}[H]
    \centering
    \includegraphics[width=0.85\textwidth]{angular.png}
    \caption{Marble-Diagramm des \texttt{first}-Operators \cite{RxJSFirstOperator}}
    \label{fig:rxjs_first}
\end{figure}


Der in Abbildung \ref{fig:rxjs_first} dargestellte \texttt{first}-Operator ist für die Lernplattform besonders dort relevant, wo nur der initiale Status einer Abfrage von Bedeutung ist. Sobald das \textit{Observable} den ersten Wert liefert, wird der Datenstrom automatisch vervollständigt, was eine effiziente Ressourcenverwaltung unterstützt.


\begin{enumerate}
\item \textbf{Asynchrone Schnittstellen-Verarbeitung:} Wenn die Lernplattform Statistiken oder neue Fragen vom Backend lädt, sorgt RxJS dafür, dass die UI nicht blockiert. Daten werden gestreamt und die Ansicht aktualisiert sich automatisch, sobald die Antwort eintrifft.
\item \textbf{Vermeidung von Inkonsistenzen:} Durch Operatoren wie \texttt{switchMap} oder \texttt{debounceTime} lassen sich komplexe Nutzerinteraktionen so steuern, dass alte, noch laufende Anfragen automatisch verworfen werden. Dies schont die Ressourcen des Clients und verhindert die Anzeige veralteter Datenstände.
\end{enumerate}


\textbf{Rendering-Strategie: Change Detection vs. Virtual DOM}


Angular verzichtet bewusst auf einen Virtual DOM wie React. Stattdessen prüft Angular direkt im echten DOM, ob sich gebundene Daten geändert haben (Change Detection). Zone.js ``monkey-patcht'' asynchrone Browser-APIs (Click, HTTP, Timer), um Change Detection automatisch auszulösen: Dabei überschreibt Zone.js zur Laufzeit native APIs wie \texttt{setTimeout}, Promises oder DOM-Events, sodass Angular bei deren Ausführung präzise Change Detection triggert, ohne manuelle Entwickler-Eingriffe \cite{web:3}.


Seit Angular 16/17 bieten \textit{Signals} eine feinere Alternative: Sie benachrichtigen \textbf{genau} die abhängigen Komponenten bei Wertänderungen, ohne Zone.js oder gesamten Komponenten-Baumlauf. Dies konvergiert Angular Richtung SolidJS-Performance, wo ebenfalls keine Virtual-DOM-Abstraktion nötig ist \cite{AngularSignalsGuide}.
\vspace{0.7cm}
\textbf{Performance-Analyse: Change Detection und Rendering}


Das Framework überwacht mithilfe von \texttt{zone.js} den Zustand der Anwendung und erkennt automatisch, wann sich Daten geändert haben, die ein Re-Rendering erfordern \cite{AngularHistory}. Dieses auf Events basierende Konzept gilt als einer der Hauptgründe für den relativ hohen Rechenaufwand im Vergleich zu moderneren Ansätzen etwa in React oder SolidJS.


In einer komplexen Lernumgebung mit vielen interaktiven Elementen (z.\,B. Wissensstand-Indikatoren) kann dies dennoch zu Performance-Einbußen führen:
\begin{itemize}
\item \textbf{Rechenintensität:} Standardmäßig prüft Angular den gesamten Komponentenbaum. Je mehr Funktionen die Plattform erhält, desto mehr CPU-Zyklen werden verbraucht.
\item \textbf{Optimierung durch OnPush:} Um eine flüssige Bedienung auf schwächeren Endgeräten sicherzustellen, kann auf die \texttt{OnPush}-Strategie umgestellt werden. Diese bewirkt, dass Komponenten nur neu gerendert werden, wenn sich Eingabeparameter (\texttt{@Input}) verändern oder ein Observable einen neuen Wert liefert. Beispiel: Ein Dashboard mit Nutzungsstatistiken aktualisiert nur jenen Teil der Ansicht, dessen Observable-Datenstrom neue Werte liefert, nicht jedoch den gesamten Baum. Diese Strategie erfordert konsequent unveränderbare Datenstrukturen (Objekt-Immutabilität), was den Lernaufwand für das Team erhöht.
\end{itemize}


\vspace{0.7cm}
\textbf{Entwickler-Produktivität: Angular CLI und Material}


Wie bei Quarkus im Backend spielt auch im Frontend das Tooling eine entscheidende Rolle. Das \textit{Angular CLI} automatisiert das Scaffolding (Erstellen von Komponenten, Services) und stellt sicher, dass alle Teammitglieder dieselbe Projektstruktur verwenden.


Für die visuelle Umsetzung bietet \textit{Angular Material} eine tief integrierte Komponentenbibliothek. Während dies die Entwicklung von Standard-Oberflächen (Buttons, Navigation) stark beschleunigt, erschwert es individuelle Design-Anpassungen: Das zugrundeliegende Design-System von Material nutzt eine komplexe Kaskade von CSS-Regeln. Dadurch muss bei tiefgreifenden Anpassungen oft mit hoher CSS-Priorität (etwa durch Überschreiben von Theme-Variablen oder Nutzung von \texttt{::ng-deep}) gearbeitet werden, was die Wartbarkeit erschwert.


\vspace{0.7cm}
\textbf{Laufzeit-Overhead: AOT-Compilation vs. „Bundle“-Größe}


Ein kritischer Punkt ist das resultierende „Bundle“, also die Dateigröße, die der Browser beim ersten Aufruf laden muss. Da Angular ein „Vollsortiment“ ist, bringt es viel Code für Funktionen mit, die eventuell gar nicht überall benötigt werden.


Der \textit{Ahead-of-Time (AOT) Compiler} kompiliert dabei TypeScript- und HTML-Templates bereits während des Build-Prozesses in effizienten JavaScript-Code. Dadurch entfällt die Auswertung der Templates zur Laufzeit, was die Startzeit reduziert und Sicherheitslücken durch dynamische Template-Ausführung minimiert \cite{AngularNG2}. Dennoch bleibt Angular im Vergleich zu leichtgewichtigeren Alternativen wie Svelte oder Vue ressourcenintensiver, was sich besonders bei mobilen Geräten bemerkbar machen kann \cite{AngMigration}.


\vspace{0.7cm}
\textbf{Fazit}


Angular bietet für ein Projekt mit zwei Teams eine hohe Stabilität durch seine klaren Vorgaben. Die „Batteries Included“-Philosophie minimiert Entscheidungsaufwände bei der Wahl von Bibliotheken. Der Overhead bei der „Bundle“-Größe fällt in neueren Versionen weniger stark ins Gewicht, da moderne Build-Tools wie \texttt{esbuild} und optimierte Tree-Shaking-Mechanismen ungenutzten Code effizient entfernen. Entscheidend bleibt jedoch die Disziplin bei der Performance-Optimierung, insbesondere bei Change Detection (z.\,B. durch OnPush oder Signals), der Lazy Loading-Strategie und einer klaren Architektur bei Observables, um unnötige Re-Renders zu vermeiden \cite{AngularFrameworkPast}.

% ==================
% 2.2.2 Flutter
% ==================
\subsection{Flutter}
\setauthor{\secondauthor}

\begin{figure}[H]
\centering
\fbox{

\includegraphics[width=0.45\textwidth]{flutter-logo.png}
}
\caption{Logo von Flutter}
\label{fig:flutter}
\end{figure}

\textbf{Historischer Kontext: Vom Experiment zur plattformübergreifenden Revolution}

Um zu begreifen, warum Flutter heute so erfolgreich ist, muss man die schwierigen Rahmenbedingungen verstehen, unter denen Entwickler Mitte der 2010er Jahre litten. Entwickler mussten sich zwischen performanten, aber teuren nativen Apps (Objective-C/Java) oder plattformübergreifenden Web-Wrappern (Cordova/PhoneGap) entscheiden, die jedoch oft träge reagierten und keine hochwertige User Experience boten \cite{FlutterHistory}.

Google stellte Flutter erstmals 2015 unter dem Codenamen "Sky" auf der Dart Developer Summit vor. Die ursprüngliche Motivation war radikal: Man wollte testen, ob eine Rendering-Engine 120 Bilder pro Sekunde erreichen kann, wenn man den Ballast der traditionellen Web-Technologien (wie das DOM-Modell) komplett abwirft. Flutter entwickelte sich schnell von einem internen Google-Experiment zu einem ernsthaften Herausforderer für Frameworks wie React Native. Während Letzteres weiterhin auf eine „Bridge“(Brücke) angewiesen war, um JavaScript-Befehle in native System-Aufrufe zu übersetzen – was oft zu Performance-Flaschenhälsen führte –, implementierte das Flutter-Team eine eigene Grafik-Engine \cite{FlutterSky}.

Mit der Veröffentlichung von Flutter 1.0 im Jahr 2018 vollzog Google den Schritt zum universellen UI-Toolkit. Die Vision weitete sich: Flutter sollte nicht mehr nur mobile Apps bedienen, sondern als "Portable UI Framework" überall dort funktionieren, wo Pixel gezeichnet werden können. Dieser Weg gipfelte 2021 in der stabilen Unterstützung für das Web. Für die Lernplattform bedeutet dies den Einsatz einer Technologie, die nicht aus der Notwendigkeit heraus entstand, Webseiten dynamischer zu machen (wie Angular), sondern mit dem Ziel, die performanteste grafische Darstellung über alle Systemgrenzen hinweg zu standardisieren.

\vspace{0.7cm}
\textbf{Technologischer Ansatz: Everything is a Widget}

Flutter bricht mit der traditionellen Web-Entwicklung, die auf HTML-Tags und CSS-Regeln basiert. Stattdessen folgt es einem rein deklarativen Ansatz, bei dem die gesamte Benutzeroberfläche aus verschachtelten Objekten, sogenannten \textit{Widgets}, besteht. Für die Lernplattform bedeutet dies eine radikale Vereinfachung des UI-Codes: Ein Button, eine Statistik-Karte oder ein ganzer Fragenkatalog sind strukturell identisch aufgebaut. Diese Konsistenz erlaubt es, Layouts extrem schnell zu prototypisieren, da Flutter eine riesige Bibliothek an vorgefertigten Material Design- und Cupertino-Widgets mitliefert, die pixelgenau auf allen Endgeräten identisch aussehen.

\vspace{0.7cm}
\textbf{Rendering: Die Skia-Engine im Web}

Der entscheidende technische Unterschied zu Frameworks wie Angular liegt im Rendering-Prozess. Während Angular das DOM (Document Object Model) des Browsers manipuliert, nutzt Flutter im Web standardmäßig das \textit{CanvasKit}, welches auf der Skia-Engine basiert.

\begin{figure}[H]
\centering
\includegraphics[width=0.85\textwidth]{FlutterRender.png}
\caption{Schichtenmodell der Flutter-Architektur \cite{FlutterArchitecture}}
\label{fig:flutter_architecture}
\end{figure}

Abbildung \ref{fig:flutter_architecture} verdeutlicht die hierarchische Struktur von Flutter: Auf oberster Ebene befinden sich die \textit{Material} und \textit{Cupertino} Widgets, die auf dem darunterliegenden \textit{Rendering Stack} (Animation, Painting, Gestures) aufbauen. Da Flutter die gesamte Benutzeroberfläche eigenständig auf ein HTML-Canvas-Element zeichnet, anstatt native DOM-Elemente zu verwenden, wird eine konsistente Performance und visuelle Treue über verschiedene Plattformen hinweg gewährleistet.
\vspace{0.7cm}

\textbf{Flutter Web: JavaScript vs. WebAssembly-Renderer}

\textit{WebAssembly (WASM)} stellt eine moderne Alternative zum klassischen JavaScript dar: kompilierten Maschinencode, der direkt im Browser läuft und deutlich höhere Performance als interpretiertes JavaScript bietet. Flutter Web unterstützt zwei Rendering-Modi:

\begin{itemize}
  \item \textbf{CanvasKit (Standard):}  
  In diesem Modus wird der Dart Code mithilfe von \texttt{dart2js} zu JavaScript kompiliert. 
  Die Kommunikation mit der Skia Rendering Engine erfolgt dabei über eine JavaScript WASM Brücke. 
  Dadurch ist CanvasKit mit allen gängigen Browsern einschließlich iOS kompatibel. 
  Allerdings entsteht durch diese Brücke ein gewisser Performance Engpass. 
  Das resultierende Bundle hat in der Regel eine Größe von etwa~1{,}5\,MB.

  \item \textbf{skwasm (WASM Modus):}  
  In diesem Modus wird der Dart Code direkt zu WebAssembly kompiliert (\texttt{dart2wasm}), 
  wodurch die Rendering Logik in einem separaten Web Worker ausgeführt werden kann. 
  Da die Kommunikation vollständig im WASM Bereich stattfindet und keine JavaScript Brücke erforderlich ist, 
  erreicht diese Variante in der Regel eine zwei- bis dreifach höhere Geschwindigkeit. 
  Allerdings setzt sie die Unterstützung von \textit{WasmGC} voraus (aktuell in Chrome~119+, Firefox~120+, 
  jedoch noch nicht auf iOS verfügbar) und erfordert zusätzliche Server Header, 
  um korrekt ausgeführt zu werden~\cite{FlutterCanvasKit}.
\end{itemize}

Bei \texttt{flutter build web --wasm} erkennt Flutter automatisch, ob der Browser skwasm unterstützt – sonst Fallback auf CanvasKit. Für datenintensive Lernplattformen mit Animationen (z. B. Echtzeit-Statistiken) ist der WASM-Modus spürbar flüssiger \cite{FlutterWasm}.

\vspace{0.7cm}
\textbf{Zustandsverwaltung: BLoC und Provider im Vergleich}

Um die Interaktivität der Lernmodi (z. B. das Umschalten zwischen Übung und Statistik) zu steuern, bietet Flutter verschiedene Muster für das State-Management an. Besonders relevant für die Lernplattform sind der \textit{Provider} und das \textit{BLoC}-Muster (Business Logic Component).

\begin{itemize}
\item \textbf{BLoC (Business Logic Component):} Trennt die UI strikt von der Logik mittels \textit{Streams}. Ein Klick auf eine Antwort sendet ein Event an den BLoC, dieser verarbeitet die Antwort (z. B. Validierung gegen das Backend) und gibt einen neuen Status an die UI zurück. Dies bietet eine exzellente Testbarkeit, erhöht aber die Komplexität des Codes durch viel "Boilerplate"-Code massiv.
\item \textbf{Provider:} Ein leichtgewichtigerer Ansatz, der Daten über den Komponentenbaum hinweg zur Verfügung stellt. Für die schnelle Implementierung von Materialverwaltungen ist dies oft effizienter, stößt aber bei sehr komplexen, ineinander verschachtelten Zuständen an seine Grenzen.
\end{itemize}

\vspace{0.7cm}
\textbf{Entwicklerproduktivität: Hot Reload und Dart}

Ein wesentlicher Faktor für die Umsetzungsgeschwindigkeit der Diplomarbeit ist das \textit{Hot Reload}-Feature. Änderungen am UI-Code werden in Millisekunden in der laufenden Anwendung sichtbar, ohne dass der aktuelle Lernfortschritt oder die Position in einem Fragenkatalog verloren geht. Dies verkürzt die Iterationszyklen im Vergleich zu Angular drastisch.

Die zugrunde liegende Sprache \textit{Dart} bietet zudem eine moderne Syntax, die Konzepte von Java und JavaScript vereint. Die starke Typisierung hilft dabei, die Datenmodelle des Backends (z. B. Fragen-Objekte) sicher im Frontend abzubilden, auch wenn die Tooling-Landschaft für Dart im Web-Bereich noch nicht die Reife von TypeScript erreicht hat.

\vspace{0.7cm}
\textbf{Schnittstellen und Web-Integration}

Obwohl Flutter als "`Multi-Platform"' -Toolkit beworben wird, zeigt sich in der Praxis eine Optimierung auf mobile Endgeräte. Der Zugriff auf Web-spezifische Funktionen (wie fortgeschrittene Browser-APIs oder die Integration spezieller JavaScript-Bibliotheken für komplexe Diagramme) ist oft umständlicher als bei nativen Web-Frameworks. Für die Lernplattform bedeutet dies, dass bei der Einbindung von Drittanbieter-Tools (z. B. für PDF-Viewer oder spezielle mathematische Formeln) oft aufwendige Platform Channels oder unfertige Community-Plugins zurückgegriffen werden muss.

\vspace{0.7cm}
\textbf{Fazit}

Flutter glänzt durch seine enorme Entwicklungsgeschwindigkeit bei der UI-Erstellung und die garantierte visuelle Konsistenz. Das "Everything is a Widget"-Konzept ist intuitiv und das Hot-Reload-Feature spart wertvolle Zeit in der Implementierungsphase. Der WASM-Modus (skwasm) ermöglicht außergewöhnliche Laufzeit-Performance, besonders für Animationen und datenintensive Anwendungen. Jedoch erkauft man sich diese Vorteile im Web-Kontext durch eine schlechtere initiale Performance (Bundle-Größe) und eine schwächere Integration in das native Web-Ökosystem. Für eine Anwendung, deren Fokus primär auf einer schnellen, interaktiven Bedienung liegt und die eventuell später als mobile App veröffentlicht werden soll, bietet Flutter eine starke Basis – sofern die Nachteile beim initialen Laden akzeptabel sind.


% ==================
% 2.2.3 React
% ==================
\subsection{React}
\setauthor{\secondauthor}

\begin{figure}[H]
\centering
\fbox{

\includegraphics[width=0.45\textwidth]{react-logo.png}
}
\caption{Logo von React}
\label{fig:react}
\end{figure}

\textbf{Historischer Kontext: Paradigmenwechsel durch Facebook}\cite{ReactHistory}\cite{ReactJSX}


Historischer Kontext: Paradigmenwechsel durch Facebook React wurde 2011
intern bei Facebook entwickelt, um die massiven Herausforderungen bei der Syn-
chronisierung von Daten in der komplexen Facebook-Timeline zu bewältigen. Der
eigentliche Auslöser war eine „Zustands-Hölle“ innerhalb der Facebook-Chat- und
Benachrichtigungs-Funktion: Oft zeigte das System eine ungelesene Nachricht an, ob-
wohl diese bereits im Chat-Fenster geöffnet war. Die damaligen Ansätze des Imperative
Programming und das Two-Way-Binding machten es fast unmöglich, den Status der UI
über verschiedene komplexe Widgets hinweg konsistent zu halten.
Jordan Walke, ein Softwareingenieur bei Facebook, entwickelte daraufhin FaxJS (den
Vorläufer von React), beeinflusst von funktionalen Programmierkonzepten. Das Ziel war
es, die UI einfach als eine Projektion des aktuellen Zustands zu betrachten: Ändert sich
der Zustand, wird die UI „einfach neu gerendert“. Zu einer Zeit, als die meisten Frame-
works (wie das frühe AngularJS) versuchten, Logik und Darstellung durch komplexes
Two-Way-Binding zu verbinden, führte React 2013 auf der JSConf US das Konzept
des „One-Way Data Flow“ und den „Virtual DOM“ ein. Die anfängliche Skepsis der
Entwicklergemeinde gegenüber JSX (HTML in JavaScript) wich schnell der Erkenntnis,
dass dieser Paradigmenwechsel – weg von der direkten Manipulation des Browsers hin
zu einem rein zustandsgesteuerten UI-Modell – die moderne Webentwicklung effizienter
und vorhersehbarer macht. Für die Entwicklung von LearnHub bedeutet der Einsatz
von React den Zugriff auf eine Technologie, die über ein Jahrzehnt hinweg für extrem da-
tenintensive Anwendungen optimiert wurde und die größte Community im Web-Bereich
vorweisen kann.

\vspace{0.7cm}

\textbf{Philosophie: Minimalismus und "Library first"}

Im Gegensatz zu Angular oder Flutter ist React eine Bibliothek, die sich primär auf die Darstellungsschicht konzentriert. Dieser minimalistische Ansatz bietet maximale Freiheit bei der Wahl zusätzlicher Werkzeuge. Für die Lernplattform bedeutet dies jedoch einen höheren Entscheidungsaufwand: Da zwei Teams arbeiten, müssen Konventionen explizit vereinbart werden.

\vspace{0.7cm}
\textbf{Technologischer Kern: Der Virtual DOM}

Der entscheidende Performance-Vorteil von React liegt im \textit{Virtual DOM}. Anstatt bei jeder Datenänderung den gesamten Baum des echten DOMs zu manipulieren, erstellt das Framework eine virtuelle Kopie im Arbeitsspeicher.

\begin{figure}[H]
\centering
\includegraphics[width=0.85\textwidth]{react_vdom_diffing.png}
\caption{Vergleich zwischen Real DOM und Virtual DOM beim Diffing-Prozess \cite{GeeksForGeeksDiffing}}
\label{fig:react_vdom_diffing}
\end{figure}

Abbildung \ref{fig:react_vdom_diffing} illustriert die Effizienz dieses Verfahrens: Wenn sich der Zustand einer Komponente ändert, wird ein neuer Virtual DOM-Baum erzeugt und mit dem vorherigen abgeglichen (\textit{Diffing}). Nur die tatsächlich veränderten Knoten (in der Abbildung markiert) werden im realen DOM der Browser-Oberfläche aktualisiert. Dieser Prozess, auch \textit{Reconciliation} genannt, verhindert unnötige Rechenoperationen und sorgt für die flüssige Darstellung der komplexen Fragen-Interfaces in LearnHub \cite{GeeksForGeeksDiffing}.

\vspace{0.7cm}
\textbf{React Fiber und Concurrent Rendering}

Seit React 16 verwendet der \textit{Fiber Reconciler} einen unterbrechbaren Diffing-Prozess. Prioritäre Updates (User-Input) werden bevorzugt, weniger dringende (Animationen) pausiert. Dies ermöglicht \textit{Concurrent Mode} und \textit{Suspense} – entscheidend für flüssige Lernplattform-Interaktionen \cite{ReactFiber}.

\vspace{0.7cm}
\textbf{React Compiler (React 19+): Automatische Optimierung}

React 19 führt den \textit{React Compiler} ein, der Code statisch analysiert und automatisch \texttt{useMemo}/\texttt{useCallback} einfügt. Dies reduziert manuelle Optimierungsarbeit und konvergiert React hin zu fein-granularen Updates – ähnlich Signals in Angular, jedoch build-time statt runtime \cite{ReactCompiler}.

\vspace{0.7cm}
\textbf{Zustandsmanagement: Hooks und Ein-Weg-Datenfluss}

React nutzt striktes "Unidirectional Data Flow": Daten fließen von Eltern zu Kind-Komponenten.

\begin{itemize}
\item \textbf{Hooks (\texttt{useState}, \texttt{useEffect}):} Diese ermöglichen es, Logik und Zustand
direkt in funktionalen Komponenten zu kapseln. Für die Lernplattform erlaubt
dies eine sehr modulare Bauweise – so kann beispielsweise die Zeitmessung eines
Prüfungsmodus als eigene Logik-Einheit (Custom Hook) geschrieben und überall
wiederverwendet werden.
\item \textbf{Prop-Drilling:} Bei tief verschachtelten Oberflächen (z. B. Das-
hboard -> Kursliste -> Themenpool -> Einzelfrage) müssen Daten oft durch
viele Ebenen gereicht werden. Um dies zu vermeiden, ist der Einsatz des Context-
API oder externer State-Manager wie Redux oder Zustand notwendig, was die
Architektur komplexer macht. \cite{ReactHooks}.
\end{itemize}
Diese ermöglichen es, Logik und Zustand
direkt in funktionalen Komponenten zu kapseln. Für die Lernplattform erlaubt
dies eine sehr modulare Bauweise – so kann beispielsweise die Zeitmessung eines
Prüfungsmodus als eigene Logik-Einheit (Custom Hook) geschrieben und überall
wiederverwendet werden.
\vspace{0.7cm}
\textbf{Entwicklerproduktivität: JSX und Ökosystem}

\textit{JSX} koppelt HTML direkt an JavaScript – lesbar, aber bricht MVC-Trennung. Das Ökosystem ist unschlagbar: Recharts für Statistiken, MUI/Tailwind für UI-Kits. Risiko: Abhängigkeits-Hölle durch viele externe Pakete.

\vspace{0.7cm}
\textbf{Rendering-Performance: Die Re-Rendering-Falle}

Ohne \texttt{React.memo}/\texttt{useMemo} rendert React oft unnötig neu. Bei LearnHubs vielen Fortschrittsbalken/Statistiken kann dies Lag verursachen. Der React Compiler löst dies weitgehend automatisch.

\vspace{0.7cm}
\textbf{Fazit}

React überzeugt durch Virtual DOM, Fiber und Compiler – die schnellste Laufzeit-Performance. Die Library-Philosophie bietet Flexibilität, erfordert aber Team-Disziplin. Für datenintensive Lernplattformen unschlagbar, besonders mit React 19s automatischer Optimierung \cite{ReactEcosystem}.


% ==================
% 2.2.4 Fazit und finale Technologieentscheidung
% ==================

\subsection{Fazit und finale Technologieentscheidung}
\setauthor{\secondauthor}

Die Gegenüberstellung der Frontend-Technologien zeigt, dass jedes Framework spezifische Schwerpunkte setzt, die die Entwicklung der Lernplattform unterschiedlich beeinflussen würden.

\begin{itemize}
\item \textbf{React} bietet zwar eine exzellente Rendering-Performance und maximale Flexibilität, erfordert jedoch aufgrund der fehlenden internen Standards (Library-Ansatz) einen hohen Abstimmungsaufwand zwischen den Teams. Das Risiko, in eine „Abhängigkeits-Hölle"' durch zu viele externe Bibliotheken zu geraten, wurde als kritisch für die Wartbarkeit eingestuft.

\item \textbf{Flutter} besticht durch seine grafische Konsistenz und die schnelle UI-Entwicklung. Die Nachteile im Web-Kontext – insbesondere die hohen initialen Ladezeiten durch CanvasKit (1,5 MB WebAssembly-Engine) und die eingeschränkte SEO-Tauglichkeit, da Suchmaschinen nur ein leeres Canvas-Element statt lesbaren HTML-Inhalten finden – machen es jedoch zu einer weniger optimalen Wahl für eine primär browserbasierte Lernplattform.

\item \textbf{Angular} erweist sich als die geeignetste Lösung für dieses Projekt. Die „Batteries Included"-Philosophie bietet eine stabile und vorstrukturierte Basis, die besonders bei der parallelen Arbeit an zwei Teilprojekten für die nötige Code-Konsistenz sorgt. Die tiefe Integration von TypeScript und RxJS ermöglicht zudem eine robuste Handhabung komplexer asynchroner Datenströme, wie sie für die interaktiven Lernmodi erforderlich sind.
\end{itemize}

Neben den rein technischen Aspekten spielte die Praxiserfahrung des Projektteams eine entscheidende Rolle bei der Wahl. Da Angular bereits fester Bestandteil des Lehrplans im Unterrichtsfach „Softwareentwicklung"' an der HTL Leonding war, sind fundierte Vorkenntnisse in der Architektur und den Konzepten des Frameworks vorhanden.

Diese Vertrautheit mit dem Werkzeug minimiert die Einarbeitungszeit erheblich und ermöglicht es dem Team, sich von Beginn an auf die Umsetzung der komplexen Logik und der User Experience zu konzentrieren, anstatt grundlegende Framework-Konzepte neu erlernen zu müssen. Die Entscheidung für Angular ist somit eine strategische Wahl, um die begrenzte Zeit der Diplomarbeit effizient zu nutzen und gleichzeitig ein professionelles, typsicheres und skalierbares Endprodukt zu garantieren.